\documentclass[pdftex,english,oribibl]{llncs}

%% Spracheinstellungen laden
\usepackage[english]{babel}

%% Schriftart in der Ausgabe/Eingabe
\usepackage[T1]{fontenc}
\usepackage{textcomp}
\usepackage[latin1]{inputenc}

%% Zitate
\usepackage[numbers]{natbib}
\bibliographystyle{abbrvnat}
%\bibliographystyle{dinat}
%\bibliographystyle{plainnat}
%\bibliographystyle{splncs}
%% Similar to option "sectionbib" but \refname instead of \bibname
\makeatletter
\renewcommand\bibsection{\section*{\refname\@mkboth{\MakeUppercase{\refname}}{\MakeUppercase{\refname}}}}
\makeatother

%% Index
%\usepackage{makeidx}
%\makeindex

%% PDF Einstellungen
% muss nach natbib geladen werden!
\usepackage{nameref}
\usepackage{varioref}
\usepackage[pdfusetitle,pdftex,colorlinks]{hyperref}
\hypersetup{pdfborder={0 0 0}}
\hypersetup{bookmarksdepth=3}
\hypersetup{bookmarksopen=true}
\hypersetup{bookmarksopenlevel=1}
\hypersetup{bookmarksnumbered=true}
\usepackage{color}
\hypersetup{colorlinks=false}

%\usepackage[section]{tocbibind}

\makeatletter
\gdef\@keywords{}
\def\keywords#1{\gdef\@keywords{#1}}
\gdef\@subtitle{}
\def\subtitle#1{\gdef\@subtitle{#1}}

%% modified from llncs
\renewenvironment{abstract}{%
  \list{}{\advance\topsep by0.35cm\relax\small%
          \leftmargin=1cm%
          \labelwidth=\z@%
          \listparindent=\z@%
          \itemindent\listparindent%
          \rightmargin\leftmargin}%
          \item[\hskip\labelsep\bfseries\abstractname]}{%
  \if!\@keywords!\else{\item[~]\item[\hskip\labelsep\bfseries\keywordname]\@keywords}\fi%
  \endlist}

\AtBeginDocument{%
  \if!\@subtitle!\else\hypersetup{pdfsubject={\@subtitle}}\fi
  \if!\@keywords!\else\hypersetup{pdfkeywords={\@keywords}}\fi
}
\makeatother

% llncs hyperref fix
\makeatletter
\providecommand*{\toclevel@author}{0}
\providecommand*{\toclevel@title}{0}
\makeatother

%% Grafiken
\usepackage[pdftex]{graphicx}
\DeclareGraphicsExtensions{.pdf,.jpg,.png}
\usepackage{subfigure}

%% Mathe
\usepackage{amsmath}
\usepackage{amssymb}

%% Listings
\usepackage{listings}
\lstset{escapechar=\%, frame=tb, basicstyle=\footnotesize}

%% Sonstiges
\newcommand{\TODO}[1]{\par\textcolor{red}{#1}\marginpar{\textcolor{red}{TODO}}}
\newcommand{\TODOX}[1]{\textcolor{red}{#1}\marginpar{\textcolor{red}{TODO}}}
\pagestyle{plain}

% Keine "Schusterjungen"
\clubpenalty = 10000
% Keine "Hurenkinder"
\widowpenalty = 10000 \displaywidowpenalty = 10000

%%%%%%%%%%%%%%%%%%%%%%%%%%%%%%%%%%%%%%%%%%%%%%%%%%%%%%%%%%%%%%%%%%%%%%%%%%%%%%%
%%% BEGIN DOCUMENT
%%%%%%%%%%%%%%%%%%%%%%%%%%%%%%%%%%%%%%%%%%%%%%%%%%%%%%%%%%%%%%%%%%%%%%%%%%%%%%%
\title{Simulation-based resilience prediction of microservice architectures}
% \subtitle{My (optional) Subtitle}
\author{Samuel Beck, Johannes G�nth�r, Christoph Zorn}
\institute{University of Stuttgart\\Reliable Software Systems (RSS)\\70569 Stuttgart, Germany}


\begin{document}

\maketitle

\begin{abstract}
	Current software simulators are tailored towards one specific purpose of conservative software simulation. Given the success of these tools it would be useful to run these tools on microservice architectures. This paper will focus on the development of a simulator that can be used for microservice architectures.
\end{abstract}

\section{Introduction}\label{sec:introduction}

	Hier werden wir beschreiben was wir erreicht haben und wie wir unser Projekt angegangen sind. Es wird au�erdem erl�utert wie wir in diesem Paper vorgehen und was f�r Themengebiete genauer betrachtet werden.
	
	%This is my introduction. %\citet{Shaw2003WritingGoodSoftwareEngineeringResearchPapersMinitutorial} wrote a paper with %hints on how to write good software engineering research papers. By the way, this was an %example for using the \textit{natbib} command \texttt{\textbackslash{}citet\{\}}.
	
	%Section~\ref{sec:anotherSection} presents everything one must know. The conclusions follow %in Section~\ref{sec:conclusions}.

%\section{Another Section}\label{sec:anotherSection}
	
	%\textit{AspectJ} can be used to weave cross-cutting concerns into Java programs %\citep{AspectJ2007}. By the way, this was an example for using the \textit{natbib} command %\texttt{\textbackslash{}citep\{\}}.
	
	%We will now demonstrate how to use subfigures (see Figure~\ref{fig:subfig}). %Figure~\ref{fig:circle} shows a circle. A star is displayed in Figure~\ref{fig:star}.
	
	%\begin{figure}
		% \centering
		%\subfigure[This is a circle.]{\label{fig:circle}
		 % \includegraphics[width=0.3\textwidth]{figures/template_circle.pdf}
		%}
		%\subfigure[This is a star.]{\label{fig:star}
		%  \includegraphics[width=0.3\textwidth]{figures/template_star.pdf}
		%}
		%\caption{A circle (a) and a star (b). Note that any caption ends with a full stop %character.}
		% \label{fig:subfig}
	%\end{figure}
  
\section{Research}\label{sec:research}

\subsection{Tools in Comparison}\label{subsec:toolsInComparison}

\subsubsection{Spigo}
	
	During our research on existing tools for microservice simulation we discovered a tool called spigo. It was written by Adrian Cockcroft a Amazon Web Services employee in the programming language go. Therefore the name spigo comes from Simulate Protocol Interactions in Go. \newline On the first view the tool looked very promising. Spigo contains a fairly simple JSON input and the structure of the parameters is intuitive. Each microservice architecture consists of multiple microservices. Each microservice has a name, a package inheritance a counter of the instances and dependencies to other microservices. The reason we like this tool is because it can simulate the occourence of an error. Spigo uses the error monkeys from the simian army. But here lies already on of the biggest disadvanteges. One can only simulate the failure of a single microservice during the execution. 

\subsubsection{Simulizar}
	
	Simulizar is an extension of the Palladio Component Model. The tool is able to adapt to changes that the software under simulation does at runtime. This helps i.e. at observing on load-balacing. 
	
\begin{itemize}
	\item analyizing self-adaptive systems (cloud-computing)
	\item non functional property prediction (performance/reliability/maintenance/cost)
	\item modeldriven architecture
	\item helps to find bottlenecks/load-/scalingproblems
\end{itemize}

\section{Simulatorspecification}\label{sec:specifictation}

\section{Simulatordocumentation}\label{sec:documentation}


\section{Conclusions}\label{sec:conclusions}

  These are my conclusions.

\bibliography{template}

\end{document}
